\documentclass[11pt, a4paper]{article}
\usepackage[a4paper, margin=2cm]{geometry}
\usepackage{fontspec}
\usepackage[english, russian, bidi=basic, provide=*]{babel}
\babelprovide[import, onchar=ids fonts]{russian}
\babelfont{rm}{Noto Sans}
\usepackage{amsmath, amssymb, amsthm}
\usepackage{booktabs}
\usepackage{enumitem}
\usepackage{hyperref}

\title{\textbf{Полные решения заданий}\\ \large Задание 5. Сингулярное разложение (SVD), тензоры и приложения линейной алгебры}
\author{Линейная алгебра}
\date{2024}

\begin{document}

\maketitle

% ==========================================
\section{Раздел 1. Сингулярное разложение (SVD)}
% ==========================================

\subsection*{Задание 1.1}
\textbf{Условие:} Найдите сингулярные числа матрицы $A = \begin{pmatrix} 3 & 0 \\ 0 & -4 \end{pmatrix}$.

\textbf{Решение:}
1. Сингулярные числа матрицы $A$ равны квадратным корням из собственных значений матрицы $A^T A$ (или модулям собственных значений матрицы $A$, если она диагональна).
2. Вычислим $A^T A$:
$$ A^T A = \begin{pmatrix} 3 & 0 \\ 0 & -4 \end{pmatrix} \begin{pmatrix} 3 & 0 \\ 0 & -4 \end{pmatrix} = \begin{pmatrix} 9 & 0 \\ 0 & 16 \end{pmatrix} $$
3. Собственные значения $A^T A$ находятся на главной диагонали: $\lambda_1 = 16, \lambda_2 = 9$.
4. Извлекая квадратные корни и упорядочивая по невозрастанию, получаем сингулярные числа:
$$ \sigma_1 = \sqrt{16} = 4, \quad \sigma_2 = \sqrt{9} = 3. $$

\textbf{Ответ:} $\sigma_1 = 4$, $\sigma_2 = 3$.

---

\subsection*{Задание 1.2}
\textbf{Условие:} Постройте сингулярное разложение $A = U \Sigma V^T$ для матрицы $A = \begin{pmatrix} 1 & 1 \\ 0 & 1 \end{pmatrix}$.

\textbf{Решение:}
1. Вычислим матрицу $A^T A$:
$$ A^T A = \begin{pmatrix} 1 & 0 \\ 1 & 1 \end{pmatrix} \begin{pmatrix} 1 & 1 \\ 0 & 1 \end{pmatrix} = \begin{pmatrix} 1 & 1 \\ 1 & 2 \end{pmatrix} $$
2. Найдем собственные значения $A^T A$:
$$ \det(A^T A - \lambda I) = \begin{vmatrix} 1-\lambda & 1 \\ 1 & 2-\lambda \end{vmatrix} = (1-\lambda)(2-\lambda) - 1 = \lambda^2 - 3\lambda + 1 = 0 $$
$$ \lambda = \frac{3 \pm \sqrt{5}}{2} \implies \lambda_1 = \frac{3+\sqrt{5}}{2} \approx 2.618, \quad \lambda_2 = \frac{3-\sqrt{5}}{2} \approx 0.382. $$
Сингулярные числа:
$$ \sigma_1 = \sqrt{\frac{3+\sqrt{5}}{2}} = \frac{1+\sqrt{5}}{2} \approx 1.618, \quad \sigma_2 = \sqrt{\frac{3-\sqrt{5}}{2}} = \frac{\sqrt{5}-1}{2} \approx 0.618. $$
Матрица $\Sigma = \begin{pmatrix} \sigma_1 & 0 \\ 0 & \sigma_2 \end{pmatrix}$.

3. Правые сингулярные векторы ($V$) --- нормированные собственные векторы $A^T A$:
Для $\lambda_1$: $(1-\lambda_1) v_1 + v_2 = 0 \implies v^{(1)} \propto \left(1, \; \frac{\sqrt{5}+1}{2}\right)^T$.
Нормируя, получаем правые сингулярные векторы $v_1, v_2$, формирующие столбцы $V$.

4. Левые сингулярные векторы ($U$) вычисляются по формуле $u_i = \frac{1}{\sigma_i} A v_i$.

\textbf{Ответ:} $A = U \Sigma V^T$, где $\Sigma = \begin{pmatrix} \frac{1+\sqrt{5}}{2} & 0 \\ 0 & \frac{\sqrt{5}-1}{2} \end{pmatrix}$, а матрицы $U$ и $V$ ортогональны.

\newpage

% ==========================================
\section{Раздел 2. Псевдообратные матрицы}
% ==========================================

\subsection*{Задание 2.1}
\textbf{Условие:} Найдите псевдообратную матрицу Мура-Пенроуза $A^+$ для вектор-столбца $A = \begin{pmatrix} 2 \\ 1 \end{pmatrix}$.

\textbf{Решение:}
Для матрицы $A$ полного столбцового ранга формула псевдообратной матрицы имеет вид:
$$ A^+ = (A^T A)^{-1} A^T $$
1. Вычислим $A^T A$:
$$ A^T A = \begin{pmatrix} 2 & 1 \end{pmatrix} \begin{pmatrix} 2 \\ 1 \end{pmatrix} = 2^2 + 1^2 = 5. $$
2. Вычислим обратное значение скаляра: $(A^T A)^{-1} = \frac{1}{5}$.
3. Находим $A^+$:
$$ A^+ = \frac{1}{5} \begin{pmatrix} 2 & 1 \end{pmatrix} = \begin{pmatrix} 0.4 & 0.2 \end{pmatrix}. $$

\textbf{Проверка свойств Мура-Пенроуза:}
$$ A A^+ A = \begin{pmatrix} 2 \\ 1 \end{pmatrix} \begin{pmatrix} 0.4 & 0.2 \end{pmatrix} \begin{pmatrix} 2 \\ 1 \end{pmatrix} = \begin{pmatrix} 0.8 & 0.4 \\ 0.4 & 0.2 \end{pmatrix} \begin{pmatrix} 2 \\ 1 \end{pmatrix} = \begin{pmatrix} 2 \\ 1 \end{pmatrix} = A. $$

\textbf{Ответ:} $A^+ = \begin{pmatrix} 0.4 & 0.2 \end{pmatrix}$.

\newpage

% ==========================================
\section{Раздел 3. Тензоры и свёртка}
% ==========================================

\subsection*{Задание 3.1}
\textbf{Условие:} Дан трехмерный тензор $T \in \mathbb{R}^{2 \times 2 \times 2}$ с элементами $T_{ijk} = i + j + k$ (где $i,j,k \in \{1, 2\}$). Вычислите результат свёртки этого тензора по первому и второму индексам: $S_k = \sum_{i=1}^2 T_{iik}$.

\textbf{Решение:}
1. Вычислим значения элементов тензора $T_{iik}$ для $i \in \{1, 2\}$ и $k \in \{1, 2\}$:
\begin{itemize}
\item Для $k = 1$:
$$ T_{111} = 1 + 1 + 1 = 3, \quad T_{221} = 2 + 2 + 1 = 5 $$
$$ S_1 = T_{111} + T_{221} = 3 + 5 = 8. $$

\item Для $k = 2$:
$$ T_{112} = 1 + 1 + 2 = 4, \quad T_{222} = 2 + 2 + 2 = 6 $$
$$ S_2 = T_{112} + T_{222} = 4 + 6 = 10. $$
\end{itemize}

2. Результатом свёртки является вектор $S = (8, 10)^T$.

\textbf{Ответ:} $S = (8, 10)^T$.

\newpage

% ==========================================
\section{Раздел 4. Бонусные задания и Теория}
% ==========================================

\subsection*{Бонус 1}
\textbf{Условие:} Докажите, что для любой матрицы $A$ псевдообратная матрица $A^+$ удовлетворяет условию $(A^+)^+ = A$.

\textbf{Доказательство:}
Псевдообратная матрица Мура-Пенроуза $X = A^+$ по определению единственным образом удовлетворяет четырём уравнениям:
1. $A X A = A$
2. $X A X = X$
3. $(A X)^T = A X$
4. $(X A)^T = X A$

Обозначим $Y = X^+ = (A^+)^+$. Матрица $Y$ должна удовлетворять соответствующим четырём уравнениям относительно $X$:
1. $X Y X = X$
2. $Y X Y = Y$
3. $(X Y)^T = X Y$
4. $(Y X)^T = Y X$

Подставим $Y = A$ в эту систему:
1. $X A X = X$ (выполняется, так как это свойство 2 для $X$)
2. $A X A = A$ (выполняется, так как это свойство 1 для $X$)
3. $(X A)^T = X A$ (выполняется, так как это свойство 4 для $X$)
4. $(A X)^T = A X$ (выполняется, так как это свойство 3 для $X$)

Поскольку матрица $A$ удовлетворяет всем четырем уравнениям для $(A^+)^+$, а псевдообратная матрица единственна, то $(A^+)^+ = A$. \qed

---

\subsection*{Теория}

\textbf{a) В чём заключается геометрический смысл сингулярного разложения (SVD)?}
\begin{proof}[Ответ]
Геометрически любое линейное преобразование, задаваемое матрицей $A$, раскладывается на три последовательных действия:
1. Поворот/отражение исходного пространства (ортогональная матрица $V^T$).
2. Растяжение/сжатие вдоль координатных осей на коэффициенты, равные сингулярным числам $\sigma_i$ (диагональная матрица $\Sigma$).
3. Поворот/отражение в пространстве образов (ортогональная матрица $U$).
Это означает, что любое линейное отображение переводит единичную сферу в эллипсоид, полуоси которого равны сингулярным числам $\sigma_i$.
\end{proof}

\textbf{b) Чем отличается ранг тензора от ранга матрицы?}
\begin{proof}[Ответ]
Ранг матрицы равен максимальному числу ее линейно независимых строк (или столбцов) и легко вычисляется полиномиальными алгоритмами (например, методом Гаусса). 

Ранг тензора (CP-ранг) определяется как минимальное число тензоров ранга 1, сумма которых дает данный тензор. В отличие от матриц:
1. Вычисление ранга тензора порядка $\ge 3$ является NP-трудной задачей.
2. Ранг тензора может превышать его максимальную размерность по любой из осей.
\end{proof}

\end{document}